\begin{table}[t]

    \begin{threeparttable}
    \centering
    \caption{Clasificación de compuestos inorgánicos}
    \label{clas_comps}
    {\setstretch{1.25}
    
        \begin{tabular}{|p{0.175\textwidth}|p{0.225\textwidth}|p{0.2\textwidth}|p{0.3\textwidth}|} \hline
            \bfseries{Nomenclatura} & \bfseries{Unión de:} & \bfseries{Tipo de compuesto} & \bfseries{Ejemplo} \\ \hline
            
            Hidrácidos & $H + \text{ No Metal}$ & Binario & $HCl$: Ácido clohrídrico \\ \hline
            
            Hidróxidos & $\text{Metal} + OH^{-}$ & Ternario & $NaOH$: Hidróxido de sodio \\ \hline
            
            Oxácidos & $H + \text{No Metal} + O$ & Ternario & $HNO_3$: Ácido nítrico \\ \hline
            
            Oxisales & $\text{Metal} + \text{Anión} + O$ & Ternario & $CaCO_3$: Carbonato de calcio \\ \hline
            
            Óxidos metálicos & $\text{Metal} + O$ & Binario & $CaO$: Óxido de calcio \\ \hline
            
            Óxidos NO metálicos & $\text{No Metal} + O$ & Binario & $H_{2}O$ \\ \hline
            
            Sales binarias & $\text{Metal} + \text{No metal}$ & Binario & $NaCl$: Cloruro de sodio \\ \hline
            
            Sales ácidas & $\text{Metal} + H + \text{Anión}$ & Poli-atómico & $KaHSO_3$: Bisulfato de potasio \\ \hline
            
        \end{tabular}
    
    }
    
    \begin{tablenotes}
        \item[] 
    \end{tablenotes}
    
    \end{threeparttable}
\end{table}

